\section*{4 Applying a Single Restriction}
In this section, we prove that the expectation of a PARITY $\circ$ AND formula is preserved under a suitable pseudorandom restriction. The cost of the restriction is only $O(\log n)$ truly random bits, the error is $\exp(-\widetilde\Omega(\log n))$ (near-optimal), and the restriction assigns values to a constant fraction of the inputs.
\subsection*{4.1 Restriction Construction}
Set $C=500$, $\bar C=2000C$, $c=1.1$, and $\beta=0.95$, and consider the following two distributions.
\begin{itemize}
\item Let $Y$ be a $\delta^3$-biased distribution over $\{0,1\}^n$ for $\delta=\min\{n^{-12\bar C},\frac12n^{-5c/(c-1)-1}\}=n^{-12,000,000}$.\footnote{No attempt was made to optimize the constants.}
\item Let $Z$ be a $\gamma$-almost $k$-wise independent distribution over $\{0,1\}^n$ with marginals $p=1-2^{-C}$, for $k=6\log n$ and $\gamma=n^{-9}$.
\end{itemize}
Our restriction is $\operatorname{Res}(Y,Z)$, i.e., $Z$ indicates where to put ? and $Y$ fills in the non-? bits.
\begin{authorclaim}{Lemma9}
Let $f$ be a read-once PARITY $\circ$ AND formula over $n$ variables of width at most $C\log n$. Then, $\operatorname{Res}(Y,Z)$ preserves the expectation of $f$ to within error $2^{-C\log n/\log\log n}$, i.e.,
\[
|\mathbb E[f(Y+Z\wedge U)]-\mathbb E[f]|\le2^{-C\log n/\log\log n}.
\]
\end{authorclaim}
\subsection*{4.2 Buckets}
Toward proving Lemma9, we first set some preliminary notations. Recall that $f$ is of the form $f=\bigoplus_{i=1}^mf_i=\bigoplus_{i=1}^m\bigwedge_{j=1}^{w_i}\ell_{ij}$, where every literal $\ell_{ij}$ is either some variable in $\{x_1,\ldots,x_n\}$ or its negation.
Set $Q=\lceil\log_c(C\log n)\rceil=O(\log\log n)$. We partition the terms of $f$ into $Q$ buckets according to their width. Namely, for each $q\in[Q]$ we define the interval $I_q=[c^{q-1},c^q)$ and define $B_q\subseteq[m]$ to be the set of indices $i$ such that $w_i\in I_q$. Also, for $q\in[Q]$ we define $F_q=\bigoplus_{i\in B_q}f_i$, so $f=\bigoplus_{q=1}^QF_q$. For every $q\in[Q]$ we further denote $m_q=|B_q|$.
We divide into two cases (Section4.3 and Section4.4) depending on whether there exists a bucket with substantially many terms. Lemma9 will follow immediately from Lemma10 and Lemma15, which cover these two cases respectively.
\subsection*{4.3 Case I -- There Exists a Heavy Bucket}
Say that bucket $q\in[Q]$ is heavy if both $m_q>3^{c^q}$ and $m_q>\log^C n$. The first case is that there exists a heavy bucket (i.e., there are many terms of roughly the same width, even relative to $q$). In this case, we will argue that $f$ itself is balanced and also that it stays balanced, w.h.p., after a pseudorandom restriction.
\begin{authorclaim}{Lemma10}
Let $f$ be a read-once PARITY $\circ$ AND formula over $n$ variables of width at most $C\log n$. Suppose there exists a heavy bucket as defined above. Then, with probability at least $1-1/n$ over $(y,z)\sim Y\times Z$,
\[
|\mathbb E[f|_{\operatorname{Res}(y,z)}]-\mathbb E[f]|\le1/n.
\]
\end{authorclaim}
Toward proving Lemma10, let us define a few more auxiliary notations. Write $f=f_{\mathrm{rest}}\oplus F_q$, where $q$ is a heavy bucket.
\begin{authorclaim}{Claim11}
It holds that $|\mathbb E[\mathfrak f]|\le1/(4n)$.
\end{authorclaim}
\begin{proof}
By the read-once property and the fact that $f_{\mathrm{rest}}$ is bounded,
\[
|\mathbb E[\mathfrak f]|=|\mathbb E[\mathfrak f_{\mathrm{rest}}]\mathbb E[\mathfrak F_q]|
\le|\mathbb E[\mathfrak F_q]|=\prod_{i\in B_q}|\mathbb E[\mathfrak f_i]|.
\]
Each term in $F_q$ has width at least $c^{q-1}$, so
\[
|\mathbb E[\mathfrak f]|\le(1-2\cdot2^{-c^{q-1}})^{m_q}\le e^{-2\cdot2^{-c^{q-1}}\cdot m_q}.
\]
Recalling that $m_q\ge3^{c^q}$, we have $2^{-c^{q-1}}\ge m_q^{-\gamma}$ for $\gamma=\log_32^{c^{-1}}<3/4$. Thus, using that fact that $m_q\ge\log^C n$,
\[
|\mathbb E[\mathfrak f]|\le e^{-2m_q^{1-\gamma}}\le e^{-2\log^{(1-\gamma)C}n}\le2^{-\log^{100}n}.
\]
\end{proof}
Next, we must analyze the bias of $f$ after the pseudorandom restriction. Let $n_q$ be the number of variables read by $F_q$. Let $b=\lceil\log_3n_q\rceil$. We will group the terms of $F_q$ into blocks, each of which reads roughly $b$ variables. To define this grouping, first observe that $b\ge\log_3m_q$, as each term reads at least one variable. Recalling that $c^q<\log_3m_q$, we know that $b>c^q$. Therefore, since each term in $F_q$ has width at most $c^q$, we can write $F_q=\bigoplus_{i=1}^Bg_i$, where each block $g_i$ reads $b_i$ variables for $b_i\in[b-\frac12c^q,b+\frac12c^q]$.
Let us now estimate $B$, the number of blocks. Since $b>c^q$, $b_i\in[b/2,3b/2]$. Also, $m_q>\log^C n$ so $b>(C/2)\log\log n$. Thus, on the one hand, $B\ge2n_q/(3b)\ge2\cdot3^b/(9b)$, and on the other hand, $B\le n_q\le3^b$.
Toward arguing that $f$ is balanced after pseudorandom restrictions, we wish to show that with high probability, $z\sim Z$ keeps many variables in many terms alive.
\begin{authorclaim}{Definition12}
For $z\in\{0,1\}^n$ and a formula $f$, we say $f$ is good under $z$ if $z$ assigns1 to at least a $(1-\beta)$-fraction of the variables $f$ reads.
\end{authorclaim}
\begin{authorclaim}{Claim13}
For a fixed $z\in\{0,1\}^n$, let $X_z\subseteq[B]$ be the set of blocks $g_i$ that are not good under $z$. Then, with probability at least $1-1/(2n)$ over $z\sim Z$, $|X_z|\le\lceil4\log n/b\rceil$.
\end{authorclaim}
\begin{proof}
Set $k_0=\lceil4\log n/b\rceil$. Let $S\subseteq[B]$ be some subset of cardinality $k_0$. We first bound the probability $p$ that every block $g_i$ for $i\in S$ is bad under $z\sim Z$. For a truly random $z\sim\operatorname{Bernoulli}(1-2^{-C})^{\otimes n}$, the above probability is bounded by
\[
\prod_{i\in S}\binom{b_i}{\beta b_i}2^{-C\beta b_i}
\le\prod_{i\in S}2^{b_i}2^{-5b_i}
\le(2^{-4\cdot b/2})^{|S|}\le n^{-8}.
\]
Now, for every $i\in[B]$, $k\ge k_0b_i$ so for $z\sim Z$, we get that $p\le n^{-8}+\gamma\le2n^{-8}$. Thus, by the union bound, with probability at most
\[
\binom B{k_0}p\le2B^{k_0}n^{-8}
\le2(3^b)^{4\log n/b}n^{-8}
\le2n^{-(2-\log3)4}\le n^{-4/3}<\frac1{2n}
\]
there will be some $S$ whose all blocks are bad. Taking the contrapositive, we infer that with probability at least $1-1/(2n)$ over $z\sim Z$, at most $k_0$ of the $g_i$-s are bad under $z$.
\end{proof}
\begin{authorclaim}{Lemma14}
With probability at least $1-1/n$ over $(y,z)\sim Y\times Z$, it holds that $|\mathbb E[\mathfrak f|_{\operatorname{Res}(y,z)}]|\le1/(2n)$.
\end{authorclaim}
\begin{proof}
Fix a good $z$, for which at most $4\log n/b$ of the $g_i$-s are not good under it. By Claim13, $z$ is good with probability at least $1-1/(2n)$. Let $B^{\mathrm{alive}}=[B]\setminus X_z$, so
\[
f=f_{\mathrm{rest}}\oplus\left(\bigoplus_{i\in B^{\mathrm{alive}}}g_i\right)\oplus\left(\bigoplus_{i\in[B]\setminus B^{\mathrm{alive}}}g_i\right).
\]
For every $i\in B^{\mathrm{alive}}$, set the following notations.
\begin{itemize}
\item For every $y\in\{0,1\}^n$, let $g_i^y$ denote the function $g_i|_{\operatorname{Res}(y,z)}$.
\item Let $I_i^{\mathrm{dead}}\subseteq[n]$ be the literals read by $g_i$ for which $z=0$. As $i\in B^{\mathrm{alive}}$, $|I_i^{\mathrm{dead}}|\le\beta b_i\le(3\beta/2)b$. Note that each literal $j\in I_i^{\mathrm{dead}}$ is set by $y\sim Y$.
\item Let $I_i^{\mathrm{alive}}\subseteq[n]$ be the literals read by $g_i$ for which $z=1$. As $i\in B^{\mathrm{alive}}$, $|I_i^{\mathrm{alive}}|\ge(1-\beta)b_i\ge((1-\beta)/2)b$.
\end{itemize}
Define the function $h_i$ so that $h_i(y)=1$ if $g_i^y$ is a nonconstant function, and0 otherwise. Namely, $h_i(y)=\bigwedge_{j\in I_i^{\mathrm{dead}}}y'_j$, where $y'_j$ is either $y_j$ or $\neg y_j$ depending on whether $y_j$ appears positively or negatively in $g_i$. Also, define $S(y)=\sum_{i\in B^{\mathrm{alive}}}h_i(y)$, where the sum is over the reals. Denote $\mu=\mathbb E[S(U)]=\sum_{i\in B^{\mathrm{alive}}}2^{-|I_i^{\mathrm{dead}}|}$, and note that $\mu\ge|B^{\mathrm{alive}}|\cdot2^{-3\beta b/2}$.
Set $\Delta S=S-\mu$. The spectral norm of the AND function is1, and so by the sub-additivity we get that $L_1(\Delta S)\le2|B^{\mathrm{alive}}|$. Set $\ell=2\lceil C\log n/(2\log(2|B^{\mathrm{alive}}|))\rceil$. By the sub-multiplicativity of the spectral norm we have that $L_1(\Delta S^\ell)\le(2|B^{\mathrm{alive}}|)^\ell\le n^C$.
For $\epsilon=1/2$, note that $\delta\le(\epsilon/2)\cdot L_1(\Delta S^\ell)^{-1}$. By Claim5, $Y$ $\epsilon$-fools the function $\Delta S^\ell$, so
\begin{equation}
|\mathbb E[(S(Y)-\mu)^\ell]-\mathbb E[(S(U)-\mu)^\ell]|\le\epsilon.\tag{5}
\end{equation}
Next, observe that $\Delta S(U)$ is the sum of zero-mean independent random variables, as the $h_i$-s are supported over disjoint set of variables. Set $A=|B^{\mathrm{alive}}|\cdot2^{-4\beta b}$. By the Chernoff bound,
\begin{align*}
\mathbb E[\Delta S(U)^\ell]
&\le\mathbb E[\Delta S(U)^\ell\mid\Delta S(U)^\ell\le A^\ell]\\
&\quad+\mathbb E[\Delta S(U)^\ell\mid\Delta S(U)^\ell\ge A^\ell]\cdot\Pr[\Delta S(U)^\ell\ge A^\ell]\\
&\le A^\ell+|B^{\mathrm{alive}}|^\ell\cdot\Pr[\Delta S(U)\ge A]\\
&\le|B^{\mathrm{alive}}|^\ell\cdot(2^{-4\beta b\ell}+e^{-2A^2/|B^{\mathrm{alive}}|}).
\end{align*}
Recall that $b>(C/2)\log\log n$, so $3^b\ge36\log n$ for a large enough $n$, and since $B\ge(2/(9b))3^b$ we get that $B\ge8\log n/b$ and $|B^{\mathrm{alive}}|\ge B-4\log n/b\ge B/2$. Next, we observe that
\[
\frac{2A^2}{|B^{\mathrm{alive}}|}=2|B^{\mathrm{alive}}|2^{-8\beta b}
\ge B\cdot2^{-8\beta b}
\ge\frac2{9b}2^{(\log3-8\beta)b}\ge2^b.
\]
As $b\ell\le Cb\log n/\log B\le C\log n$, we can conclude that $2^b\ge4\beta b\ell$ and so $e^{-2A^2/|B^{\mathrm{alive}}|}\le2^{-4\beta b\ell}$, which implies that $\mathbb E[\Delta S(U)^\ell]\le2|B^{\mathrm{alive}}|^\ell\cdot2^{-4\beta b\ell}$.
Using Equation(5) and the above bound yields a bound on $\mathbb E[\Delta S(Y)^\ell]$. By Markov's inequality,
\begin{equation}
\Pr[S(Y)<\mu/2]\le\frac{\mathbb E[(S(Y)-\mu)^\ell]}{(\mu/2)^\ell}
\le\frac{\epsilon+2|B^{\mathrm{alive}}|^\ell\cdot2^{-4\beta b\ell}}{(\mu/2)^\ell}
\le\left(\frac{8|B^{\mathrm{alive}}|2^{-4\beta b}}{\mu}\right)^\ell.\tag{6}
\end{equation}
Recalling that $\mu\ge|B^{\mathrm{alive}}|\cdot2^{-3\beta b/2}$, Equation(6) becomes
\[
\Pr[S(Y)<\mu/2]\le(8\cdot2^{(-4\beta+3\beta/2)b})^\ell
<2^{-2\beta b\ell}\le2^{-\beta C\log n/2}\le\frac1{2n},
\]
where we have used the fact that $b\ell\ge C\log n/4$.
Overall, with probability at least $1-1/(2n)$ over $y\sim Y$, $g_i^y$ is nonconstant for at least $\mu/2$ of the $i$-s, and recall that each such $g_i^y$ is over at least $(1-\beta)b_i$ variables. Fix such a good $y$, and let $\mathcal G\subseteq[B^{\mathrm{alive}}]$ be the set of nonconstant $g_i^y$-s. Again, we can write
\[
\bigoplus_{i\in B^{\mathrm{alive}}}g_i^y
=\left(\bigoplus_{i\in\mathcal G}g_i^y\right)\oplus\left(\bigoplus_{i\in B^{\mathrm{alive}}\setminus\mathcal G}g_i^y\right)
\mathrel{\mathop:}=t_1\oplus t_2.
\]
Similarly to Claim11, in order to bound the bias of $f|_{\operatorname{Res}(Y,Z)}$ it is sufficient to bound the bias of $t_1$, and so
\[
|\mathbb E[\mathfrak t_1]|\le(1-2^{-3b/2})^{\mu/2}.
\]
Using the fact that $\mu\ge(B/2)\cdot2^{-3\beta b/2}\ge(1/(9b))2^{(\log3-3\beta/2)b}>2^{301b/200}$, we get
\[
|\mathbb E[\mathfrak t_1]|\le e^{-2^{-3b/2}2^{301b/200}}\le e^{-\log^{C/400}n}\le\frac1{2n}.
\]
\end{proof}
\begin{proof}[Proof of Lemma10]
Finally, the fact that with probability at least $1-1/n$ over $(y,z)\sim Y\times Z$, $|\widetilde f_z(y)-\mathbb E[\mathfrak f]|\le1/n$, follows immediately from Claim11 and Lemma14.
\end{proof}
